\begingroup
\setlength{\LTcapwidth}{\textwidth}
\renewcommand{\arraystretch}{1.15}
\begin{longtable}{@{}p{4.4cm}p{4.0cm}p{4.0cm}@{}}
\caption{Ringkasan karakteristik kanal glukosa pada dataset OhioT1DM.}
\label{tab:tab3.3_karakteristik_dataset} \\
\toprule
\textbf{Karakteristik} & \textbf{CGM} & \textbf{Finger-stick} \\
\midrule
\endfirsthead

\multicolumn{3}{@{}l}{\small\emph{Tabel \thetable{} (lanjutan)}} \\
\toprule
\textbf{Karakteristik} & \textbf{CGM} & \textbf{Finger-stick} \\
\midrule
\endhead

\bottomrule
\endlastfoot

Jumlah baris & 166.533 & 4.566 \\

Pembacaan per pasien & 13.878 & 380 \\

Median selang antar-pembacaan & 5,0 menit & 124,8 menit \\

Rerata glukosa (SD) & 159,6 (60,7) mg/dL & 164,9 (69,6) mg/dL \\

Hipoglikemia $<70$ mg/dL & 3,28\% & 4,58\% \\

Dalam rentang target 70--180 mg/dL & 63,42\% & 58,34\% \\

Hiperglikemia $>180$ mg/dL & 33,30\% & 37,08\% \\

Split train/test & 134.790 / 31.743 & 3.801 / 765 \\
\end{longtable}
\endgroup
